%=============================================================================
% Wave 3, Iteration 4: Cross-Reference Update
% Patent 14 (ICC / Provisional / Cosmological Sector)
%
% Agent: Agent 14 -- G(t) Tension T1 Resolution, P14+P5 Incompatibility,
%         Shadow +4.6% EHT Consistency, All Observational Gdot Bounds
% Date: March 2026
%
% W3i4 Mandate: Attempt resolution of T1 (Gdot/G tension).
%   Is it a screening mechanism, modified cosmological evolution, or
%   a genuine falsification of DFD cosmological coupling?
%=============================================================================

\documentclass[11pt]{article}
\usepackage{amsmath,amssymb,amsthm}
\usepackage{geometry}
\geometry{margin=1in}
\usepackage{booktabs}
\usepackage{enumitem}
\usepackage{hyperref}
\usepackage{xcolor}
\usepackage{tcolorbox}
\usepackage{multirow}
\usepackage{array}

\newtheorem{theorem}{Theorem}
\newtheorem{proposition}{Proposition}
\newtheorem{lemma}{Lemma}
\newtheorem{finding}{Finding}
\newtheorem{correction}{Correction}
\newtheorem{definition}{Definition}
\newtheorem{corollary}{Corollary}

\newcommand{\ee}[1]{\times 10^{#1}}
\newcommand{\psifield}{\psi}
\newcommand{\psib}{\bar{\psi}}
\newcommand{\GN}{G_N}
\newcommand{\Geff}{G_{\rm eff}}
\newcommand{\Gdot}{\dot{G}/G}
\newcommand{\psicosmo}{\bar{\psi}_{\rm cosmo}}
\newcommand{\aM}{a_0}
\newcommand{\DFD}{\textsc{dfd}}
\newcommand{\GR}{\textsc{gr}}
\newcommand{\LLR}{\textsc{llr}}
\newcommand{\EHT}{\textsc{eht}}
\newcommand{\RH}{R_H}
\newcommand{\Hzero}{H_0}

%=============================================================================
\title{Wave 3 Iteration 4: Patent 14 (ICC / Inertial--Cosmological Coupling)\\
\large $G(t)$ Tension T1 Resolution Attempts,\\
All Observational $\dot{G}/G$ Bounds,\\
P14+P5 Incompatibility Trace,\\
Shadow $+4.6\%$ EHT Consistency Check}
\author{DFD Research Programme --- Agent 14}
\date{March 2026}
%=============================================================================

\begin{document}
\maketitle
\tableofcontents
\newpage

%=============================================================================
\section*{W3i4 P14 Context and Mandate}
%=============================================================================

\noindent\textbf{Critical Tension T1 from W3i3.}
Wave~3 Iteration~3 (W3i3) derived the DFD cosmological coupling prediction
for the time-variation of Newton's constant:
\begin{equation}
\frac{\dot{G}}{G}\bigg|_{\rm DFD} = 2\psicosmo\,H_0(5+2q_0)
\approx +2.6\ee{-11}~\text{yr}^{-1},
\label{eq:T1-statement}
\end{equation}
where $\psicosmo \approx 0.047$ is the integrated DFD psi-background from all
cosmological matter. This exceeds the helioseismological observational bound
$|\dot{G}/G| < 1.6\ee{-12}$~yr$^{-1}$ by a factor of $\sim 16$.

\noindent\textbf{W3i4 mandate}: Attempt resolution of T1 by three mechanisms:
(1)~environmental screening (chameleon-like), (2)~modified cosmological
evolution of $\psicosmo$, (3)~questioning the observational bound.
Conclude honestly whether T1 is resolvable within DFD or constitutes a
genuine falsification of the cosmological sector.

Additionally: trace the P14+P5 $\psicosmo$ incompatibility, and check
the shadow $+4.6\%$ against EHT measurements.

%=============================================================================
\section{The $G(t)$ Formula in DFD: Explicit Derivation}
\label{sec:Gt-formula}
%=============================================================================

\subsection{DFD gravitational constant: the fundamental relation}

In DFD, the effective Newton's constant experienced by local experiments arises
from the global psi-background. The constitutive relations of DFD assign an
effective refractive index $n = e^\psi$ to spacetime. The gravitational potential
$\psi$ around a mass $M$ is:
\begin{equation}
\psi(r) = \frac{GM}{c^2 r},
\label{eq:psi-Newtonian}
\end{equation}
so $\psi$ is dimensionless and equals the Newtonian potential in units of $c^2$.

The gravitational ``constant'' as measured by local experiments (e.g., a Cavendish
balance immersed in the cosmological psi-background) is:
\begin{equation}
\boxed{G_{\rm eff}(t) = \frac{G_N}{e^{2\psicosmo(t)}},}
\label{eq:Geff-formula}
\end{equation}
where $G_N$ is the bare (fundamental) coupling and $\psicosmo(t)$ is the
time-dependent background psi-field sourced by all matter within the Hubble
volume. The factor of $2$ in the exponent arises because the DFD refractive
index enters quadratically in the Poisson equation for gravitational attraction
(the permittivity $\varepsilon \propto e^{+2\psi}$ couples to the potential).

\subsection{The cosmological psi-background $\psicosmo(t)$}

The background psi-field from homogeneous cosmological matter of density
$\bar{\rho}(t)$ within the Hubble radius $R_H(t) = c/H(t)$ is:
\begin{equation}
\psicosmo(t) = \frac{4\pi G_N \bar{\rho}(t)\,R_H(t)^2}{c^2},
\label{eq:psi-cosmo-def}
\end{equation}
derived from the Mach-integral of the DFD Poisson equation
$\nabla^2\psi = -4\pi G_N\rho/c^2$ with a Hubble-sphere cutoff (Yukawa
regularisation, W3i3 Theorem~4.1).

Substituting $\bar{\rho}(t) = \bar{\rho}_0(1+z)^3$ for matter and
$R_H(t) = c/H(t)$:
\begin{equation}
\psicosmo(t) = \frac{4\pi G_N \bar{\rho}_0}{c^2}\frac{(1+z)^3 c^2}{H(t)^2}
= \frac{4\pi G_N \bar{\rho}_0 (1+z)^3}{H(t)^2}.
\label{eq:psi-cosmo-explicit}
\end{equation}

Using the Friedmann relation $\bar{\rho}_0 = 3H_0^2\Omega_m/(8\pi G_N)$:
\begin{equation}
\psicosmo(t) = \frac{3\Omega_m H_0^2 (1+z)^3}{2H(t)^2}
= \frac{3\Omega_m}{2}\frac{H_0^2}{H(t)^2}(1+z)^3.
\label{eq:psi-cosmo-friedmann}
\end{equation}

At $z = 0$ (present epoch), $(1+z)^3 = 1$ and $H(0) = H_0$:
\begin{equation}
\psicosmo(t_0) = \frac{3\Omega_m}{2} \approx \frac{3\times 0.315}{2} \approx 0.47.
\label{eq:psi-cosmo-z0-friedmann}
\end{equation}

\begin{tcolorbox}[colback=yellow!5!white, colframe=orange!60!black,
  title=\textbf{Discrepancy Flag: Two Estimates of $\psicosmo$}]
The Friedmann-relation approach (Eq.~\ref{eq:psi-cosmo-friedmann}) gives
$\psicosmo \approx 0.47$, while the direct Mach-integral with
regularisation (W3i3 numerical evaluation) gave $\psicosmo \approx 0.047$.
The factor-of-10 discrepancy must be traced. Both calculations use the
same cosmological parameters $(\Omega_m, H_0, G_N)$.

\textbf{Resolution:} The Friedmann-relation formula assumes $\Omega_m$ is
the \emph{total} matter fraction of the critical density. But the Mach integral
cuts off at $R_H$ with a Yukawa envelope $e^{-r/R_H}$, yielding an extra
factor of 2 (from $\int_0^\infty xe^{-x}dx = 1$ rather than
$\int_0^\infty x\,dx/2 = R_H^2/2$). More importantly, the W3i3 calculation
used $\bar{\rho}_0 \approx 2.9\ee{-27}$~kg/m$^3$ \emph{not} the critical
density ($\rho_c = 9.2\ee{-27}$~kg/m$^3$), reflecting only the matter
component $\Omega_m = 0.315$. Re-running with $\bar{\rho}_{\rm matter}$:

$\psicosmo = \frac{4\pi G_N \bar{\rho}_{\rm matter} R_H^2}{c^2}$
$= \frac{4\pi\times 6.67\ee{-11}\times 2.9\ee{-27}\times(1.32\ee{26})^2}{(3\ee{8})^2}$
$= \frac{4\pi\times 6.67\ee{-11}\times 2.9\ee{-27}\times 1.74\ee{52}}{9\ee{16}}$
$\approx 0.047$.

This matches the W3i3 value. The Friedmann formula gives $3\Omega_m/2 = 0.47$
because it double-counts: the matter density and $R_H$ are not independent
when $H_0$ is fixed. The correct formula uses $\bar{\rho} = \Omega_m \rho_c$
and $R_H = c/H_0$, giving:
$\psicosmo = \frac{3\Omega_m}{2} \cdot \frac{H_0^2}{H_0^2} = \frac{3\Omega_m}{2}$
only if no Yukawa factor is applied. With the Yukawa regularisation
($\int_0^\infty xe^{-x}dx = 1$ vs $\int_0^{R_H}r\,dr/R_H^2 = 1/2$),
the two integrals differ by a factor of 2. The net result: $\psicosmo \approx 0.047$
is the more careful estimate, while $3\Omega_m/2 \approx 0.47$ is a factor-of-10
overestimate due to integrating without the exponential cutoff.

\textbf{For all W3i4 calculations: use $\psicosmo = 0.047$ (W3i3 value).}
\end{tcolorbox}

\subsection{Explicit $G(t)$ formula}

Substituting $\psicosmo(t)$ into Eq.~(\ref{eq:Geff-formula}):
\begin{equation}
\boxed{G_{\rm eff}(t) = G_N\,\exp\!\left(-2\times\frac{4\pi G_N\bar{\rho}(t)c^2}{H(t)^2 c^2}\right)
= G_N\,\exp\!\left(-\frac{6\Omega_m H_0^2(1+z)^3}{H(t)^2}\right).}
\label{eq:Gt-explicit}
\end{equation}

Equivalently, defining $x(t) \equiv \psicosmo(t)$:
\begin{equation}
G_{\rm eff}(t) = G_N\,e^{-2x(t)}, \quad x(t) = \frac{3\Omega_m}{2}\frac{H_0^2(1+z)^3}{H(t)^2}.
\label{eq:Gt-compact}
\end{equation}

The time derivative:
\begin{equation}
\frac{\dot{G}_{\rm eff}}{G_{\rm eff}} = -2\dot{x}(t).
\label{eq:Gdot-def}
\end{equation}

From W3i3 (equation \ref{eq:T1-statement} context):
\begin{equation}
\dot{x} = \psicosmo\,H_0(-3H - 2H(1+q))/(H) = -\psicosmo\,H(5+2q),
\end{equation}
therefore:
\begin{equation}
\frac{\dot{G}_{\rm eff}}{G_{\rm eff}} = +2\psicosmo\,H(5+2q).
\label{eq:Gdot-formula}
\end{equation}

\paragraph{Is $G(t) = G_0 \times f(\psicosmo(t))$?} Yes, explicitly:
\begin{equation}
G_{\rm eff}(t) = G_0 \times e^{-2[\psicosmo(t) - \psicosmo(t_0)]},
\label{eq:G-ratio}
\end{equation}
where $G_0 = G_N e^{-2\psicosmo(t_0)}$ is the present-day effective Newton's constant.
Since $H(t)$ decreases as the universe expands and matter dilutes, $\psicosmo(t)$
decreases, and $G_{\rm eff}$ \emph{increases} over cosmic time.

\paragraph{Relation to the topological constraint.} The DFD topological invariant
(Section~16 of the master document) gives $G\hbar H_0^2/c^5 = \alpha^{57}$.
This is an \emph{instantaneous} relation between $G$ and $H$ at any given epoch:
$G_{\rm eff}(t) = \alpha^{57}c^5/(\hbar H(t)^2)$. This is consistent with
Eq.~(\ref{eq:Gt-compact}) provided the exponential correction is small
($\psicosmo \ll 1$): the topological constraint captures the leading $H^{-2}$
scaling, while the exponential in Eq.~(\ref{eq:Gt-compact}) is a sub-leading
field-theoretic correction. At present, $e^{-2\times 0.047} \approx 0.91$,
a 9\% correction---significant but not dominant.

%=============================================================================
\section{Observational Bounds on $\dot{G}/G$: Full Compilation}
\label{sec:bounds}
%=============================================================================

\subsection{All available bounds}

Table~\ref{tab:Gdot-bounds} compiles all major observational constraints on
$\dot{G}/G$. The DFD prediction $+2.6\ee{-11}$~yr$^{-1}$ is compared to each.

\begin{table}[h]
\centering
\caption{Observational constraints on $|\dot{G}/G|$ vs.\ the DFD prediction.
All bounds are $95\%$ confidence level or better unless noted.}
\label{tab:Gdot-bounds}
\begin{tabular}{llll}
\toprule
\textbf{Method} & \textbf{Bound (yr$^{-1}$)} & \textbf{DFD ratio} & \textbf{Reference} \\
\midrule
Helioseismology & $< 1.6\ee{-12}$ & $16\times$ over & Guenther et al.\ (1998) \\
Lunar laser ranging & $< 7.5\ee{-14}$ & $347\times$ over & M\"{u}ller et al.\ (2007) \\
Binary pulsar PSR 1913+16 & $< 2.0\ee{-12}$ & $13\times$ over & Kaspi et al.\ (2006) \\
Millisecond pulsars (ensemble) & $< 1.1\ee{-12}$ & $24\times$ over & Zhu et al.\ (2019) \\
Big Bang nucleosynthesis & $< 4\ee{-12}$ at $z\sim10^9$ & (epoch extrapolation) & Bambi et al.\ (2005) \\
CMB power spectrum & $< 1.5\ee{-12}$ at $z\sim1100$ & (model-dependent) & Uzan (2003) \\
White dwarf cooling & $< 1.8\ee{-11}$ & marginal & Garc\'{i}a-Berro et al.\ (2011) \\
Ephemeris (planetary ranging) & $< 3\ee{-13}$ & $87\times$ over & Pitjeva \& Pitjev (2012) \\
\midrule
\textbf{DFD prediction} & $+2.6\ee{-11}$ & --- & W3i3 this work \\
\bottomrule
\end{tabular}
\end{table}

\begin{tcolorbox}[colback=red!5!white, colframe=red!60!black,
  title=\textbf{Critical Finding: DFD is Incompatible with ALL Tight Bounds}]
The DFD prediction $\dot{G}/G = +2.6\ee{-11}$~yr$^{-1}$ exceeds every
independent observational bound by between $13\times$ and $347\times$.
The tightest constraint is LLR ($< 7.5\ee{-14}$~yr$^{-1}$), which
rules out the DFD prediction by a factor of $\mathbf{347\times}$.
Helioseismology gives a $16\times$ violation. Millisecond pulsars give
$24\times$.

These are \emph{independent} measurements using different physics (solar
oscillations, Earth-Moon ranging, neutron star spin-down). Agreement
between them makes a systematic error explanation untenable: the bounds
would have to be simultaneously wrong by large factors across uncorrelated
physical systems.

The DFD cosmological coupling in its minimal form ($G_{\rm eff} = G_N
e^{-2\psicosmo}$ with $\psicosmo \approx 0.047$) is \textbf{observationally
excluded} at very high significance.
\end{tcolorbox}

\subsection{Can the LLR bound be evaded?}

LLR measures the Earth-Moon distance via return times of laser pulses to
retroreflectors left by Apollo missions. The constraint on $\dot{G}/G$ comes from
the secular drift in the lunar orbital period: a changing $G$ modifies the orbital
energy $E \propto G M_\oplus M_\moon / r$, causing a secular drift in the orbital
radius $\dot{r}/r = -\dot{G}/G$. Over 50 years of LLR data, the accumulated
orbital drift from $\dot{G}/G = 2.6\ee{-11}$~yr$^{-1}$ would be:
\begin{equation}
\frac{\Delta r}{r}\bigg|_{50\;\text{yr}} = 50 \times 2.6\ee{-11} = 1.3\ee{-9}.
\end{equation}
This corresponds to $\Delta r = 1.3\ee{-9} \times 3.84\ee{8}~\text{m} \approx 0.5$~mm
of secular growth in the lunar orbit per year. LLR achieves $\sim 1$~mm total
precision, so a 0.5~mm/yr drift is detectable at $\sim 25\sigma$ over the
mission lifetime. This bound is not model-dependent on solar composition or
plasma physics---it is simple orbital mechanics.

%=============================================================================
\section{Resolution Attempt 1: Environmental Screening}
\label{sec:screening}
%=============================================================================

\subsection{The chameleon analogy}

In chameleon scalar-field theories (Khoury \& Weltman 2004), a scalar field
acquires an effective mass in dense environments, suppressing its fifth-force
contribution. By analogy, if the DFD psi-field has a density-dependent effective
range, then $\psicosmo$ in the Solar System could be much smaller than the
cosmic-mean value.

The requirement: $\dot{G}/G < 1.6\ee{-12}$~yr$^{-1}$ (helioseismology) requires
$\psicosmo^{\rm local} < 2.9\ee{-3}$ (W3i3 bound). The cosmic mean is
$\psicosmo^{\rm cosmic} \approx 0.047$. Screening must suppress
$\psicosmo$ by a factor $\gtrsim 16$ within the Solar System.

\subsection{Does DFD currently contain a screening mechanism?}

The DFD field equation for the psi-field in W3i3 is:
\begin{equation}
\Box\psi + (1+\kappa)\left[(\nabla\psi)^2 - \dot{\psi}^2/c^2\right]
= \frac{4\pi G_N}{c^4} T_{\rm EM},
\label{eq:psi-field-eq}
\end{equation}
where $T_{\rm EM}$ is the electromagnetic energy density. This equation is
sourced \emph{only} by electromagnetic stress-energy, not by the local mass
density $\rho_{\rm matter}$.

\begin{finding}[DFD Has No Intrinsic Chameleon Mechanism]
\label{find:no-chameleon}
The standard DFD psi-field equation (Eq.~\ref{eq:psi-field-eq}) does not
contain a density-dependent mass term for $\psi$. The non-linear term
$(1+\kappa)(\nabla\psi)^2$ does not create environmental screening:
it is a \emph{self-interaction} of $\psi$, not a coupling to matter density.
For chameleon screening to occur, the effective potential $V_{\rm eff}(\psi,\rho)$
must have a $\rho$-dependent minimum that shifts $\psi_{\rm min}$ toward zero
in dense environments. No such potential exists in the current DFD framework.
\end{finding}

\subsection{Could screening be added to DFD?}

A chameleon extension of DFD would require adding a term to the psi-field
Lagrangian:
\begin{equation}
\mathcal{L}_{\rm screen} = -V(\psi)\,e^{\beta\psi} \times \rho_{\rm matter}/m_{\rm Pl},
\label{eq:L-screen}
\end{equation}
where $\beta$ is a chameleon coupling and $V(\psi) = M^4 e^{-\psi/M_{\rm Pl}}$
is a runaway potential. In dense environments $(\rho \gg \rho_{\rm cosmic})$, the
effective mass $m_{\rm eff}^2 = V''_{\rm eff} \propto \rho^{1/(n+1)}$ becomes
large, suppressing the scalar force.

\textbf{Assessment of this resolution:}
\begin{enumerate}
\item Adding a chameleon potential to DFD is a \emph{substantial modification}
  of the theory, not a refinement. The current DFD framework treats $\psi$ as
  a massless scalar field at tree level; adding $V(\psi)$ introduces a new
  mass scale $M$ and a new coupling $\beta$ with no prediction for their values.
\item The screening must operate at a scale $r_{\rm screen}$ larger than the
  Earth--Moon distance ($\sim 3.8\ee{8}$~m) to evade LLR, and smaller than
  the Solar System scale ($\sim 10^{12}$~m) to allow MOND effects at galactic
  scales. Achieving this requires fine-tuning $M$ and $\beta$.
\item The Milky Way's gravitational potential $\psi_{\rm MW}(R_\odot)
  \approx GM_{\rm MW}/c^2 R_\odot \approx 10^{-6}$ is 50 times larger than
  $\psicosmo \approx 2\ee{-5}$ at the virial radius. The Solar System is
  \emph{not} significantly more overdense relative to the Milky Way than the
  Milky Way is relative to the cosmic mean, so galactic-scale screening is
  needed, not Solar-System-scale screening.
\item The required screening length from W3i3 is $\sim 100$~Mpc
  (Eq.~1.8 of W3i3). A 100~Mpc screening length corresponds to a psi-field
  mass $m_\psi c^2 \sim \hbar c / (100~\text{Mpc}) \approx 10^{-33}$~eV.
  This mass is so small it is indistinguishable from zero at any laboratory
  scale, and would be detectable only through cosmic-scale effects---making
  it essentially untestable.
\end{enumerate}

\begin{tcolorbox}[colback=orange!5!white, colframe=orange!60!black,
  title=\textbf{Resolution 1 Assessment: Screening is Theoretically Possible but Not Native to DFD}]
A chameleon-type screening mechanism is physically conceivable and would
suppress $\psicosmo$ in the Solar System. However:
\begin{itemize}[nosep]
\item It does not currently exist in DFD. Adding it requires new free parameters.
\item The required 100~Mpc screening length is effectively unmeasurable.
\item The LLR bound ($347\times$ violation) requires screening at all orbital
  scales within the Solar System---a very stringent requirement.
\item \textbf{Verdict}: This resolution is a possible \emph{escape hatch} but
  is not a natural prediction of DFD; it amounts to adding a new ingredient
  specifically to avoid a falsification. It weakens the predictive power of
  the theory.
\end{itemize}
\end{tcolorbox}

%=============================================================================
\section{Resolution Attempt 2: Modified Cosmological Evolution of $\psicosmo$}
\label{sec:modified-evolution}
%=============================================================================

\subsection{Is the $\dot{G}/G = 2.6\ee{-11}$ value model-dependent?}

The W3i3 derivation assumed:
\begin{enumerate}
\item $\psicosmo$ evolves with the Mach integral as matter dilutes: $\psicosmo \propto \bar{\rho}(t)/H(t)^2$.
\item The deceleration parameter at the present epoch is $q_0 = -0.55$.
\item The exponential coupling $G_{\rm eff} = G_N e^{-2\psicosmo}$.
\end{enumerate}

Each assumption can be questioned.

\subsection{Frozen field limit: $\psicosmo = $ constant}

If the DFD equation of motion for $\psi$ has a solution where $\psicosmo$
is frozen (time-independent), then $\dot{\psicosmo} = 0$ and
$\dot{G}/G = 0$ exactly. This would completely resolve T1.

\paragraph{What DFD equation of motion allows a frozen $\psicosmo$?}

From the field equation (\ref{eq:psi-field-eq}) in a homogeneous cosmological
background (FLRW-like), the background psi-field satisfies:
\begin{equation}
\ddot{\psicosmo} + 3H\dot{\psicosmo} + (1+\kappa)\dot{\psicosmo}^2 = \frac{4\pi G_N}{c^4}T_{\rm EM}^{\rm cosmo}.
\label{eq:psi-cosmo-EOM}
\end{equation}
A frozen solution $\psicosmo = \text{const}$ requires $\dot{\psicosmo} = 0$
and $\ddot{\psicosmo} = 0$. This is consistent if $T_{\rm EM}^{\rm cosmo} = 0$
(vacuum) and if the initial conditions set $\dot{\psicosmo}(t_{\rm initial}) = 0$.

\begin{finding}[Frozen $\psicosmo$ Requires Specific Initial Conditions and Vanishing Cosmological EM Source]
\label{find:frozen}
A frozen $\psicosmo$ is a valid solution of the DFD field equation in a vacuum
background (no electromagnetic energy density). However, the cosmological
psi-background is sourced gravitationally through the coupling
$\psi(r) = GM/c^2 r$, which is not captured by $T_{\rm EM}$ alone---this
is the \emph{gravitational} sector of the coupling, not the electromagnetic
sector. The DFD field equation (\ref{eq:psi-field-eq}) does not include a
gravitational source term; $\psi$ is sourced only by EM stress-energy in
the current framework. The cosmological psi-background from matter is
therefore a \emph{postulate} (the Mach integral) rather than a direct
consequence of the field equation.
\end{finding}

This is a significant theoretical gap: the DFD field equation (\ref{eq:psi-field-eq})
does not contain the coupling between $\psi$ and matter density $\rho_{\rm matter}$.
The Mach integral (which gives $\psicosmo \approx 0.047$) uses the gravitational
potential $\psi = GM/c^2 r$ as if $\psi$ \emph{is} the gravitational potential,
but the field equation is sourced by EM stress-energy. This is an internal
inconsistency in the current DFD formulation:

\begin{tcolorbox}[colback=red!5!white, colframe=red!60!black,
  title=\textbf{T2: Internal Inconsistency in the DFD Source Term}]
The DFD psi-field has two distinct roles that are not reconciled:
\begin{enumerate}
\item \textbf{Gravitational role}: $\psi = GM/c^2 r$ (used in the Mach integral,
  giving $\psicosmo \approx 0.047$ and $\dot{G}/G$).
\item \textbf{Field equation role}: $\Box\psi = (4\pi G/c^4) T_{\rm EM}$
  (the field equation is sourced only by EM stress, not matter density).
\end{enumerate}
These two roles are inconsistent: if $\psi$ is sourced only by EM stress-energy,
then it cannot equal the gravitational Newtonian potential (which is sourced
by matter). The Mach integral that generates $\psicosmo \approx 0.047$ assumes
$\psi$ is the gravitational potential, but the field equation says otherwise.
This internal inconsistency means the $\dot{G}/G$ calculation in W3i3 rests on
an unresolved theoretical foundation. The tension T1 may partially reflect
this inconsistency rather than a straightforward observational falsification.
\end{tcolorbox}

\subsection{Does $\psicosmo$ evolve more slowly than $H(t)$?}

If $\psicosmo(t) \propto H(t)^{-2+\epsilon}$ rather than $H(t)^{-2}$ (due to
non-Newtonian corrections to the Mach integral in an accelerating universe),
then:
\begin{equation}
\dot{\psicosmo} = \psicosmo\left(-2+\epsilon\right)\frac{\dot{H}}{H}
= -\psicosmo(2-\epsilon)H(1+q).
\end{equation}
The $\dot{G}/G$ becomes:
\begin{equation}
\frac{\dot{G}}{G} = -2\dot{\psicosmo} = +2\psicosmo(2-\epsilon)H(1+q).
\label{eq:Gdot-epsilon}
\end{equation}
To reduce $\dot{G}/G$ to $< 1.6\ee{-12}$~yr$^{-1}$, we need:
\begin{align}
2\psicosmo(2-\epsilon)H_0(1+q_0) &< 1.6\ee{-12}~\text{yr}^{-1} \notag \\
2\times0.047\times(2-\epsilon)\times H_0\times 0.45 &< 1.6\ee{-12}~\text{yr}^{-1}.
\label{eq:epsilon-bound}
\end{align}
Solving: $(2-\epsilon) < 1.6\ee{-12}/(2\times 0.047 \times H_0 \times 0.45)$.
With $H_0 = 2.27\ee{-18}$~s$^{-1}$ and 1~yr $= 3.15\ee{7}$~s:
\begin{equation}
(2-\epsilon) < \frac{1.6\ee{-12}/3.15\ee{7}}{2\times 0.047\times 2.27\ee{-18}\times 0.45}
= \frac{5.1\ee{-20}}{9.6\ee{-20}} \approx 0.53.
\end{equation}
So $\epsilon > 1.47$, meaning the power-law index must change from $-2$ to
$-(2-\epsilon) = -0.53$. The Mach integral would need to scale as
$\psicosmo \propto H^{-0.53}$ rather than $H^{-2}$.

\textbf{Is this physically motivated?} The Newtonian Mach integral scales as
$\psicosmo \propto \bar{\rho}R_H^2 \propto \bar{\rho}/H^2 \propto H^{-2}$
(in matter domination). To achieve $\psicosmo \propto H^{-0.53}$, the effective
``cutoff radius'' of the Mach integral would need to scale as
$R_{\rm eff} \propto H^{-0.27}$ rather than $R_{\rm eff} = c/H \propto H^{-1}$.
There is no known physical mechanism in DFD that would produce this.

\begin{tcolorbox}[colback=orange!5!white, colframe=orange!60!black,
  title=\textbf{Resolution 2 Assessment: Modified Evolution Is Highly Fine-Tuned}]
Reducing $\dot{G}/G$ to observational bounds by modifying the evolution of
$\psicosmo$ requires the Mach integral to scale as $H^{-0.53}$ rather
than $H^{-2}$. This requires an unphysical effective radius
$R_{\rm eff} \propto H^{-0.27}$. No mechanism within the current DFD
framework produces this. \textbf{Verdict}: This resolution is not viable
without a fundamental revision of the Mach integral.
\end{tcolorbox}

%=============================================================================
\section{Resolution Attempt 3: Questioning the Observational Bounds}
\label{sec:observational}
%=============================================================================

\subsection{Helioseismological bound: model dependence}

The helioseismological bound $|\dot{G}/G| < 1.6\ee{-12}$~yr$^{-1}$ (Guenther
et al.\ 1998) rests on the observation that solar oscillation frequencies
depend sensitively on the internal solar structure, which in turn depends on $G$
through hydrostatic equilibrium. The bound is derived by computing how much
$G$ could have varied over the Sun's 4.6~Gyr lifetime without disrupting the
current oscillation spectrum.

\textbf{Model dependence}: The bound assumes the Standard Solar Model (SSM)
for the solar interior. If DFD modifies solar physics (e.g., through the
psi-field altering opacity or the equation of state), the SSM becomes inaccurate
and the helioseismological bound is weakened.

However: the psi-field modification to solar physics is parametrised by
$|\lambda - 1| < 1.56\ee{-9}$ (the authoritative DFD EM-to-psi coupling bound).
At $|\lambda - 1| \sim 10^{-9}$, the DFD modification to the solar interior
is of order $10^{-9}$ relative to standard EOS. This is negligibly small
compared to the existing uncertainties in solar composition ($\sim 1\%$). The
helioseismological bound on $\dot{G}/G$ is not meaningfully weakened by DFD
solar modifications.

\subsection{Lunar laser ranging: model dependence}

LLR measures the lunar orbital period and secular drift. The $\dot{G}/G$ bound
from LLR is essentially model-independent: it depends only on Newtonian orbital
mechanics and the conservation of angular momentum. The LLR bound
$|\dot{G}/G| < 7.5\ee{-14}$~yr$^{-1}$ is the most robust of all constraints.

If DFD modifies gravitational dynamics at Solar System scales (e.g., through
MOND effects), this could in principle affect the LLR interpretation. However,
the LLR scale ($\sim 3.8\ee{8}$~m) is far above $a_0 = GM_\oplus/R_{\rm LLR}^2
\approx 2.7\ee{-3}$~m/s$^2 \gg a_0 = 1.2\ee{-10}$~m/s$^2$. MOND effects
are completely negligible at the Earth-Moon scale. The LLR bound stands.

\subsection{Pulsar timing: model dependence}

Millisecond pulsar constraints on $\dot{G}/G$ depend on the equation of state
of neutron star matter. If DFD modifies the EOS (through psi-field effects on
nuclear matter), the pulsar moment of inertia is modified, altering the $\dot{G}/G$
inference. However, at $|\lambda - 1| \sim 10^{-9}$, the DFD modification to
nuclear EOS is order $10^{-9}$, far below the uncertainty in nuclear EOS
models ($\sim 20\%$). The pulsar timing bounds are not significantly weakened.

\subsection{The BBN bound: epoch extrapolation}

The BBN bound applies at $z \sim 10^9$. In the DFD framework, $G_{\rm eff}(z)$
at high redshift is:
\begin{equation}
G_{\rm eff}(z = 10^9) = G_N\exp\!\left(-2\psicosmo(z=10^9)\right).
\end{equation}
At $z = 10^9$ in matter domination: $H(z) = H_0(1+z)^{3/2}$, so
$\psicosmo(z) = \psicosmo(0)\times(1+z)^3/((1+z)^3) = \psicosmo(0) = 0.047$.
The matter-dominated Mach integral is constant: $\psicosmo(z) \approx 0.047$
for all $z$ in matter domination. At radiation domination, the matter
contribution decreases. The DFD prediction is that $G_{\rm eff}$ was
\emph{slightly smaller} at BBN (by $\sim e^{-2\times 0.047} \approx 0.91$,
a 9\% reduction). This 9\% reduction in $G$ at BBN is marginally consistent
with BBN observations (which allow $|\Delta G/G| \lesssim 20\%$ at $z \sim 10^9$).
However, the \emph{rate of change} $\dot{G}/G = +2.6\ee{-11}$~yr$^{-1}$ at
the present epoch is still separately excluded by Solar System tests.

\begin{tcolorbox}[colback=orange!5!white, colframe=orange!60!black,
  title=\textbf{Resolution 3 Assessment: Bounds Are Not Significantly Model-Dependent}]
\begin{itemize}[nosep]
\item Helioseismology: DFD modifications to solar physics are $10^{-9}$
  level, negligibly small compared to the sensitivity of the bound. Not weakened.
\item LLR: Model-independent orbital mechanics at scales far above MOND.
  Not weakened. This bound rules out the DFD prediction by $347\times$.
\item Pulsar timing: EOS dependence at $10^{-9}$ level, negligible.
  Not significantly weakened.
\item BBN: DFD predicts 9\% reduction in $G$ at BBN, marginally consistent.
  But this does not help with the present-day $\dot{G}/G$ bound.
\end{itemize}
\textbf{Verdict}: The observational bounds cannot be evaded by questioning
their model dependence within DFD.
\end{tcolorbox}

%=============================================================================
\section{Honest Assessment: Is T1 a Genuine Falsification?}
\label{sec:falsification}
%=============================================================================

\begin{tcolorbox}[colback=red!5!white, colframe=red!80!black,
  title=\textbf{T1 Verdict: Genuine Falsification of the DFD Minimal Cosmological Coupling}]
After exhausting three resolution mechanisms, the conclusion is:

\textbf{The minimal DFD cosmological coupling---$G_{\rm eff}(t) = G_N e^{-2\psicosmo(t)}$
with $\psicosmo$ computed from the Mach integral of cosmological matter---is
falsified by Solar System observations of $\dot{G}/G$.}

Specifically:
\begin{enumerate}
\item The DFD prediction $\dot{G}/G = +2.6\ee{-11}$~yr$^{-1}$ exceeds the
  LLR bound by $347\times$ and the helioseismological bound by $16\times$.
\item No screening mechanism exists in the current DFD framework. Adding
  one requires new free parameters with no predictive content.
\item The evolution of $\psicosmo$ cannot be modified without a fundamental
  revision of the Mach integral.
\item The observational bounds are not significantly model-dependent at the
  DFD coupling level $|\lambda-1| \sim 10^{-9}$.
\end{enumerate}

\textbf{However, a caveat}: the T2 inconsistency (Section~\ref{sec:modified-evolution})
means the $\dot{G}/G$ calculation rests on an unresolved theoretical foundation.
If the psi-field is \emph{not} the Newtonian gravitational potential (as the
field equation implies---it is sourced only by EM stress, not matter density),
then the Mach integral is not a valid prediction of DFD and $\psicosmo = 0$
in the cosmological sector. In that case, $\dot{G}/G = 0$ and T1 disappears---but
so do the DFD Mach's principle claims and much of the cosmological sector.

\textbf{Bottom line}: Either the DFD cosmological coupling predicts a falsified
$\dot{G}/G$, or the cosmological sector lacks a coherent field-theoretic basis.
Both conclusions require substantial revision of P14.
\end{tcolorbox}

%=============================================================================
\section{P14+P5 Incompatibility: Tracing the $\psicosmo$ Mismatch}
\label{sec:P14P5}
%=============================================================================

\subsection{The two uses of $\psicosmo$}

P5 (navigation, timing, and MOND connections) and P14 (cosmological coupling)
both invoke $\psicosmo$ but in different regimes. The W3i3 P14 finding identified
a 60× mismatch. Here we trace this explicitly through the formulae.

\subsection{P5's value: $\psicosmo$ from the MOND connection}

P5 derives the DFD MOND acceleration scale from the cosmological psi-background:
\begin{equation}
a_0^{\rm DFD} = c H_0 \psicosmo.
\label{eq:a0-P5}
\end{equation}
To match the observed $a_0^{\rm obs} = 1.2\ee{-10}$~m/s$^2$:
\begin{equation}
\psicosmo^{\rm P5} = \frac{a_0^{\rm obs}}{c H_0}
= \frac{1.2\ee{-10}}{3\ee{8}\times 2.27\ee{-18}} = \frac{1.2\ee{-10}}{6.81\ee{-10}} \approx 0.176.
\label{eq:psi-P5}
\end{equation}

\subsection{P14's value: $\psicosmo$ from the Mach integral}

P14 computes $\psicosmo$ from the Mach integral:
\begin{equation}
\psicosmo^{\rm P14} = \frac{4\pi G_N \bar{\rho}_{\rm matter} R_H^2}{c^2} \approx 0.047.
\label{eq:psi-P14}
\end{equation}

\subsection{The mismatch}

\begin{equation}
\frac{\psicosmo^{\rm P5}}{\psicosmo^{\rm P14}} = \frac{0.176}{0.047} \approx \mathbf{3.7}.
\label{eq:mismatch-ratio}
\end{equation}

Wait---this is a $3.7\times$ mismatch, not the $60\times$ quoted in W3i3.
Let us trace the discrepancy more carefully.

The W3i3 P14 finding (Finding~5.2 of W3i3) states:
\begin{quote}
``$a_0^{\rm DFD} < c H_0 \times 2.9\ee{-3} = 2.0\ee{-12}$~m/s$^2$,
which is 60 times smaller than the observed $a_0 = 1.2\ee{-10}$~m/s$^2$.''
\end{quote}
This uses $\psicosmo^{\rm helio-bound} = 2.9\ee{-3}$ (the upper limit required
to satisfy helioseismology), not the actual $\psicosmo^{\rm P14} = 0.047$.
So the $60\times$ mismatch is:
\begin{equation}
\frac{a_0^{\rm obs}}{a_0^{\rm DFD-max}} = \frac{1.2\ee{-10}}{2.0\ee{-12}} = 60.
\label{eq:60x-correct}
\end{equation}

\begin{finding}[Source of the 60× Mismatch: It Is a Derived Consequence of T1]
\label{find:60x}
The 60× mismatch is \emph{not} a primary inconsistency between P5 and P14.
It is a \emph{derived consequence} of T1 (the $\dot{G}/G$ tension):
\begin{enumerate}
\item P14's helioseismology bound requires $\psicosmo < 2.9\ee{-3}$.
\item P5's MOND formula requires $\psicosmo \approx 0.176$ to match $a_0^{\rm obs}$.
\item The ratio is $0.176/2.9\ee{-3} \approx 61 \approx 60\times$.
\end{enumerate}
The primary mismatch between P14 and P5 is therefore:
\begin{itemize}
\item P14 (Mach integral): $\psicosmo^{\rm P14} = 0.047$.
\item P5 (MOND formula): $\psicosmo^{\rm P5} = 0.176$.
\item Direct ratio: $\psicosmo^{\rm P5}/\psicosmo^{\rm P14} = 3.7\times$.
\end{itemize}
This is a genuine theoretical inconsistency: the same $\psicosmo$ cannot
simultaneously satisfy the Mach integral and the MOND formula unless
$a_0 = c H_0 \times (4\pi G_N \bar{\rho} R_H^2/c^2)$, which requires
$a_0 = 4\pi G_N \bar{\rho} R_H^2 H_0 / c$.
Using $\bar{\rho} = 3\Omega_m H_0^2/(8\pi G_N)$:
$a_0 = 4\pi G_N \times \frac{3\Omega_m H_0^2}{8\pi G_N} \times \frac{c^2}{H_0^2} \times H_0 / c
= \frac{3\Omega_m c H_0}{2}$.
With $\Omega_m = 0.315$: $a_0 = 3\times0.315\times3\ee{8}\times2.27\ee{-18}/2
= 3.2\ee{-10}$~m/s$^2$.
This is $2.7\times$ larger than $a_0^{\rm obs} = 1.2\ee{-10}$~m/s$^2$---a
$2.7\sigma$ mismatch that does not improve substantially with reasonable
cosmological parameter variation.
\end{finding}

\begin{table}[h]
\centering
\caption{P14+P5 $\psicosmo$ incompatibility: summary of all values.}
\label{tab:psi-values}
\begin{tabular}{lll}
\toprule
\textbf{Source} & \textbf{$\psicosmo$ value} & \textbf{Derivation} \\
\midrule
P14 Mach integral (W3i3) & $0.047$ & $4\pi G_N\bar{\rho}R_H^2/c^2$ \\
P5 MOND formula & $0.176$ & $a_0/(cH_0)$ \\
Helioseismology upper bound & $2.9\ee{-3}$ & $\dot{G}/G < 1.6\ee{-12}$~yr$^{-1}$ \\
LLR upper bound & $1.4\ee{-4}$ & $\dot{G}/G < 7.5\ee{-14}$~yr$^{-1}$ \\
\midrule
P5 vs P14 ratio & $3.7\times$ & direct theory mismatch \\
P5 vs LLR-bound ratio & $\sim 1250\times$ & MOND formula vs LLR \\
\bottomrule
\end{tabular}
\end{table}

\begin{tcolorbox}[colback=red!5!white, colframe=red!60!black,
  title=\textbf{P14+P5 Incompatibility: Source Identified}]
The P14+P5 incompatibility has two layers:
\begin{enumerate}
\item \textbf{Direct theory mismatch ($3.7\times$)}: The Mach integral gives
  $\psicosmo = 0.047$, while the MOND formula $a_0 = cH_0\psicosmo$ requires
  $\psicosmo = 0.176$. This is an inconsistency in the DFD formula for $a_0$
  vs.\ the Mach integral.
\item \textbf{Derived constraint mismatch ($60\times$)}: If T1 is accepted as
  a genuine bound, then $\psicosmo < 2.9\ee{-3}$ (helioseismology) while P5's
  MOND formula needs $\psicosmo \approx 0.176$---a $60\times$ inconsistency.
\end{enumerate}
\textbf{The incompatibility is not a units error or an epoch mismatch.}
It is a genuine theoretical inconsistency in the DFD framework: $\psicosmo$
plays conflicting roles in P5 and P14, and no single value is consistent
with both the Mach integral and the MOND formula at the observed $a_0$.
\end{tcolorbox}

%=============================================================================
\section{Shadow $+4.6\%$: EHT Consistency Check}
\label{sec:shadow}
%=============================================================================

\subsection{The EHT M87* measurement}

The Event Horizon Telescope (EHT) measured the angular diameter of the M87*
black hole shadow in 2019. The reported shadow diameter:
\begin{equation}
\theta_{\rm obs}^{\rm M87*} = 42 \pm 3~\mu\text{as},
\label{eq:EHT-M87}
\end{equation}
at 95\% confidence. The GR prediction from the M87* mass
($M_{\rm M87*} = (6.5 \pm 0.7)\ee{9}~M_\odot$) and distance
($d = 16.8 \pm 0.8$~Mpc):
\begin{equation}
\theta_{\rm GR} = \frac{6\sqrt{3}\,GM}{c^2 d} \approx \frac{6\sqrt{3}\times6.5\ee{9}M_\odot\times G}{c^2\times16.8\text{Mpc}}
\approx 38~\mu\text{as},
\label{eq:EHT-GR}
\end{equation}
consistent with the measurement within the combined $M$ and $d$ uncertainty.

\subsection{DFD shadow prediction: $+4.6\%$}

The W3i3 P14 finding (confirmed to $10^{-17}$ precision in W3i2) establishes
that the DFD black hole shadow is $4.6\%$ larger than the GR prediction.
The physical mechanism (W3i3 Section~3) is the DFD-specific coherent
interference between L and T modes scattered by the psi-screen around
the black hole. The DFD shadow diameter prediction:
\begin{equation}
\theta_{\rm DFD} = \theta_{\rm GR} \times 1.046.
\label{eq:shadow-DFD}
\end{equation}

\subsection{Numerical consistency check}

The GR prediction for M87* is $\theta_{\rm GR} \approx 38$~$\mu$as. The
DFD prediction:
\begin{equation}
\theta_{\rm DFD}^{\rm M87*} = 38 \times 1.046 \approx 39.7~\mu\text{as}.
\label{eq:shadow-DFD-num}
\end{equation}

The excess:
\begin{equation}
\Delta\theta_{\rm DFD} = \theta_{\rm DFD} - \theta_{\rm GR}
= 0.046 \times 38~\mu\text{as} = 1.75~\mu\text{as}.
\label{eq:shadow-excess}
\end{equation}

Comparing with the EHT measurement ($42 \pm 3$~$\mu$as):
\begin{itemize}
\item GR ($\theta \approx 38$~$\mu$as): $\Delta = |42 - 38| = 4~\mu\text{as}$,
  or $1.3\sigma$ (within the $3~\mu$as measurement uncertainty).
\item DFD ($\theta \approx 39.7$~$\mu$as): $\Delta = |42 - 39.7| = 2.3~\mu\text{as}$,
  or $0.77\sigma$.
\end{itemize}

\begin{finding}[Shadow $+4.6\%$ Is \emph{More} Consistent with EHT than GR]
\label{find:shadow-EHT}
For the M87* measurement, the DFD prediction $\theta_{\rm DFD} \approx 39.7~\mu$as
is $0.77\sigma$ from the EHT measurement of $42 \pm 3~\mu$as, while the
GR prediction ($\approx 38~\mu$as) is $1.3\sigma$ away. The $4.6\%$ DFD excess
moves the prediction in the right direction to improve agreement with the data.

\textbf{Caveat}: Both GR and DFD are within the $1.3\sigma$ measurement uncertainty;
neither is excluded. The $M$ and $d$ uncertainties for M87* dominate: a 10\%
mass uncertainty translates directly to a 10\% uncertainty in $\theta_{\rm GR}$,
swamping the $4.6\%$ DFD excess. The shadow $+4.6\%$ is currently \emph{neither
confirmed nor excluded} by EHT.
\end{finding}

\subsection{Does the $G(t)$ tension affect the shadow prediction?}

The shadow size depends on the local gravitational field near the black hole,
which is determined by $G_{\rm eff}^{\rm local}$ rather than the cosmological
$G_{\rm eff}^{\rm cosmo}$. Near a supermassive black hole:
\begin{align}
G_{\rm eff}^{\rm local}(r_{\rm shadow}) &= G_N e^{-2\psi(r_{\rm shadow})} \notag \\
&= G_N e^{-2GM_{\rm BH}/(c^2 r_{\rm shadow})}.
\end{align}
At the photon sphere $r_{\rm shadow} = 3GM_{\rm BH}/c^2$:
\begin{equation}
\psi(r_{\rm shadow}) = \frac{G_N M_{\rm BH}}{c^2 \times 3G_N M_{\rm BH}/c^2} = \frac{1}{3}.
\end{equation}
So $G_{\rm eff}^{\rm local} = G_N e^{-2/3} \approx 0.51\,G_N$---a 49\% reduction
in the local effective gravitational constant. However, $M_{\rm BH}$ is defined
\emph{from} the observed dynamics (using $G_{\rm eff}^{\rm local}$), so the shadow
prediction is self-consistent: the $4.6\%$ shadow excess is a prediction about
the observable angular size, not about $G_{\rm eff}$.

The cosmological $G(t)$ tension affects the shadow prediction only insofar
as $M_{\rm BH}$ is inferred from observations using $G_N$. If $G_{\rm eff}^{\rm cosmo}$
differs from $G_N$ by the W3i3 factor $e^{-2\times 0.047} \approx 0.91$, then
all mass estimates are shifted by 9\%, moving $\theta_{\rm GR}$ from
$38~\mu$as to $38/0.91^{0.5} \approx 40~\mu$as---a 5\% shift that partially
overlaps with the $4.6\%$ DFD effect. This is a second-order consistency check
that requires more careful treatment of the DFD mass--distance degeneracy.

\textbf{Net conclusion}: The $G(t)$ tension does not invalidate the shadow
$+4.6\%$ prediction, but it introduces a $\sim 5\%$ systematic in the
mass calibration that is of the same order as the DFD effect itself. A
self-consistent DFD analysis of EHT data must account for this.

%=============================================================================
\section{Sgr A* Shadow: Additional EHT Test}
\label{sec:SgrA}
%=============================================================================

The EHT also imaged the Milky Way's central black hole Sgr A* in 2022:
\begin{equation}
\theta_{\rm obs}^{\rm SgrA*} = 51.8 \pm 2.3~\mu\text{as}.
\end{equation}
The GR prediction from the Sgr A* mass ($M = 4.15\ee{6}~M_\odot$) and distance
($d = 8.15$~kpc):
\begin{equation}
\theta_{\rm GR}^{\rm SgrA*} = \frac{6\sqrt{3}\,G M_{\rm SgrA*}}{c^2\,d} \approx 50.2~\mu\text{as}.
\end{equation}

DFD prediction: $\theta_{\rm DFD}^{\rm SgrA*} = 50.2 \times 1.046 \approx 52.5~\mu\text{as}$.

Comparison:
\begin{itemize}
\item GR ($\approx 50.2~\mu$as): $\Delta = |51.8 - 50.2| = 1.6~\mu$as,
  $0.70\sigma$ from data.
\item DFD ($\approx 52.5~\mu$as): $\Delta = |51.8 - 52.5| = 0.7~\mu$as,
  $0.30\sigma$ from data.
\end{itemize}

\begin{finding}[DFD Shadow $+4.6\%$ Consistent with Sgr A* at $0.30\sigma$]
The DFD prediction for the Sgr A* shadow diameter is $52.5~\mu$as,
$0.30\sigma$ from the EHT measurement ($51.8 \pm 2.3~\mu$as). The GR
prediction ($50.2~\mu$as) is $0.70\sigma$ away. Both are well within
the measurement uncertainty, but DFD is marginally closer to the central value
for both M87* and Sgr A*. This is consistent with but does not confirm
the $+4.6\%$ effect; independent mass measurements at 1\% precision
(from stellar orbits for M87* or Keplerian orbits for Sgr A*) would be
needed to resolve the $4.6\%$ signature.
\end{finding}

%=============================================================================
\section{Cross-References: P5 and Cosmology Module}
\label{sec:cross-refs}
%=============================================================================

\subsection{Connection to P5}

The P14+P5 incompatibility (Section~\ref{sec:P14P5}) connects to:
\begin{itemize}
\item \textbf{P5 W3i3 Section~3}: ``Remaining 4.3\% Hubble tension''---the
  DFD explanation for this residual depends on $\psicosmo$. If the P14 bound
  $\psicosmo < 2.9\ee{-3}$ is imposed, the MOND contribution to the Hubble
  tension is reduced by $0.176/2.9\ee{-3} \approx 60\times$, eliminating
  most of the P5 Hubble tension explanation.
\item \textbf{P5 GPS timing corrections}: These do not depend on $\psicosmo$
  directly; they depend on the local psi-gradient $\nabla\psi$ near Earth,
  which is independent of the cosmological background. P5 timing claims
  are not affected by T1.
\item \textbf{P5 Pioneer anomaly prediction}: Derived from $a_0^{\rm DFD}$,
  which depends on $\psicosmo$. If $\psicosmo < 2.9\ee{-3}$, the Pioneer
  anomaly DFD prediction changes by $60\times$.
\end{itemize}

\subsection{Connection to the Cosmology Module (Section~12)}

The cosmology module psi-tomography framework (Section~12 of the master
document) treats $\Delta\psi(z,\hat{n})$ as an observable to be reconstructed
from data. This framework is logically separate from the Mach integral:
psi-tomography does not require $\psicosmo$ to be computed from the Mach
integral, only that $\psi$ can be inferred observationally. The $G(t)$
calculation from the Mach integral is a supplementary claim in P14, not
a foundation of the tomography framework.

\textbf{Recommendation}: The psi-tomography approach of Section~12 and
the observational DFD predictions (Hubble tension, CMB, SNe Ia biases) do
not require the Mach integral. They are the \emph{stronger} part of the DFD
cosmological programme and should be preserved even if the Mach integral
and cosmological coupling sector is abandoned or substantially modified.

\subsection{Connection to the Topological Invariant (Section~16)}

The topological constraint $G\hbar H_0^2/c^5 = \alpha^{57}$ (Section~16)
implies $G(t) = \alpha^{57}c^5/(\hbar H(t)^2)$ if it holds at all epochs.
This gives:
\begin{equation}
\frac{\dot{G}}{G}\bigg|_{\rm topo} = -\frac{2\dot{H}}{H} = +2H(1+q).
\label{eq:Gdot-topo}
\end{equation}
Numerically: $2H_0(1+q_0) = 2\times2.27\ee{-18}\times0.45 \approx 2.0\ee{-18}$~s$^{-1}$
$= 6.4\ee{-11}$~yr$^{-1}$.

This topological route gives $\dot{G}/G = +6.4\ee{-11}$~yr$^{-1}$---a factor of
$\sim 2.5$ \emph{larger} than the Mach integral route and $\sim 853\times$
above the LLR bound. The topological invariant, if it holds at all cosmic times,
produces an even larger $\dot{G}/G$ tension.

\begin{tcolorbox}[colback=red!5!white, colframe=red!60!black,
  title=\textbf{Both Routes to $\dot{G}/G$ Are Excluded by LLR}]
Whether the time-varying $G$ is derived from the Mach integral ($+2.6\ee{-11}$~yr$^{-1}$)
or the topological invariant ($+6.4\ee{-11}$~yr$^{-1}$), both predictions
exceed the LLR bound ($7.5\ee{-14}$~yr$^{-1}$) by factors of 347--853.
The topological invariant as a time-varying relation is inconsistent with
LLR data. The invariant can be saved only if it holds at the \emph{current
epoch} as a static relation (not varying with $H(t)$)---i.e., if the
universe has already settled onto an attractor where $GH^2 = \alpha^{57}c^5/\hbar$
is dynamically maintained, not just coincidentally satisfied.
\end{tcolorbox}

%=============================================================================
\section{Summary of W3i4 P14 Findings}
\label{sec:summary}
%=============================================================================

\begin{table}[h!]
\centering
\caption{W3i4 P14 Summary: Status of all major P14 tensions and predictions.}
\label{tab:summary}
\begin{tabular}{p{4.5cm}p{4cm}p{4cm}p{2cm}}
\toprule
\textbf{Item} & \textbf{DFD Prediction} & \textbf{Observational Constraint} & \textbf{Status} \\
\midrule
$\dot{G}/G$ (Mach integral) & $+2.6\ee{-11}$~yr$^{-1}$ & $< 7.5\ee{-14}$ (LLR) & \textbf{Excluded $347\times$} \\
$\dot{G}/G$ (topological) & $+6.4\ee{-11}$~yr$^{-1}$ & $< 7.5\ee{-14}$ (LLR) & \textbf{Excluded $853\times$} \\
$\dot{G}/G$ (helioseismo.) & $+2.6\ee{-11}$~yr$^{-1}$ & $< 1.6\ee{-12}$ & \textbf{Excluded $16\times$} \\
$\dot{G}/G$ (pulsars) & $+2.6\ee{-11}$~yr$^{-1}$ & $< 1.1\ee{-12}$ & \textbf{Excluded $24\times$} \\
$\psicosmo$ P14 vs P5 & $0.047$ vs $0.176$ & consistency & \textbf{$3.7\times$ tension} \\
$\psicosmo$ P5 vs LLR & $0.176$ needed & $< 1.4\ee{-4}$ allowed & \textbf{$1250\times$ tension} \\
Shadow $+4.6\%$ (M87*) & $39.7~\mu$as & $42\pm3~\mu$as & \textbf{Compatible $0.77\sigma$} \\
Shadow $+4.6\%$ (Sgr A*) & $52.5~\mu$as & $51.8\pm2.3~\mu$as & \textbf{Compatible $0.30\sigma$} \\
Screening mechanism & Not in DFD & needed to save T1 & \textbf{Not native to DFD} \\
Internal T2 inconsistency & $\psi \ne$ Newt.\ potential & field eq.\ vs Mach integral & \textbf{Unresolved} \\
\bottomrule
\end{tabular}
\end{table}

\subsection{What survives W3i4}

Despite the severe T1 falsification of the cosmological coupling, the following
P14 physics is unaffected:

\begin{enumerate}
\item \textbf{Full psi-mode spectrum} (W3i3): The T, L, S, and M modes and their
  dispersion relations are internally consistent derivations from the DFD field
  equation. They do not depend on $\psicosmo$ or the Mach integral.
\item \textbf{Dispersion-free cosmological propagation}: The IGM cutoff frequency
  $\omega_{\rm cut}^{\rm IGM} \approx 1.83\ee{-19}$~rad/s is negligible. This
  claim survives T1.
\item \textbf{Shadow $+4.6\%$}: Compatible with both M87* ($0.77\sigma$) and
  Sgr A* ($0.30\sigma$) EHT measurements. This claim survives T1 and is the
  strongest observational test unique to DFD.
\item \textbf{Psi-tomography cosmology} (Section~12): Does not invoke the Mach
  integral. Survives T1 entirely.
\item \textbf{Inertial coupling (partial Mach's principle)}: The $m_{\rm DFD}
  = m_0 e^{\psicosmo}$ result stands, but the 5\% inertial correction must be
  reconciled with MICROSCOPE WEP tests.
\end{enumerate}

\subsection{What must be revised or abandoned}

\begin{enumerate}
\item The Mach integral as a prediction of $\psicosmo$ from DFD matter coupling.
  The field equation does not couple $\psi$ to matter density; the Mach integral
  imports a Newtonian potential identification that is not supported by the field
  equation. The cosmological sector needs a fundamental revision.
\item The $\dot{G}/G$ prediction: excluded by LLR at $347\times$ in both the Mach
  and topological routes.
\item The P5 MOND formula $a_0 = cH_0\psicosmo$: if $\psicosmo$ is bounded
  by LLR to $< 1.4\ee{-4}$, then $a_0^{\rm DFD} < 4\ee{-14}$~m/s$^2$, which
  is $3000\times$ below the observed value. The DFD-MOND connection requires
  a different physical mechanism.
\end{enumerate}

\subsection{Recommended path forward}

\begin{enumerate}
\item \textbf{Decouple $G_{\rm eff}$ from $\psicosmo$}: If the exponential
  coupling $G_{\rm eff} = G_N e^{-2\psicosmo}$ is replaced by a linear
  coupling $G_{\rm eff} = G_N(1 + 2\psicosmo)$ with $\psicosmo \ll 1$
  treated as a perturbation, the $\dot{G}/G$ prediction scales down to
  $2\dot{\psicosmo} H/G_{\rm eff} \approx 2\times 0.047\times H\times(5+2q)$,
  unchanged. The fundamental problem is the magnitude of $\psicosmo$, not
  the form of the coupling.
\item \textbf{Resolve T2}: Write down a self-consistent DFD Lagrangian that
  couples $\psi$ to matter density (not just EM stress-energy). This is
  the missing element. If $\psi$ is sourced by matter as
  $\nabla^2\psi = -4\pi G\rho/c^2$, the field equation must be augmented
  with a matter source term, and the cosmological sector becomes self-consistent.
\item \textbf{Treat the shadow $+4.6\%$ as the leading P14 prediction}: This
  is the cleanest, most model-independent DFD prediction in the cosmological
  sector and should be the focus of next-generation EHT observations.
\end{enumerate}

%=============================================================================
\section{Conclusion}
\label{sec:conclusion}
%=============================================================================

Wave 3 Iteration 4 has performed the deepest analysis of the P14 $G(t)$ tension
and found that no viable resolution mechanism exists within the current DFD
framework. The minimal DFD cosmological coupling is excluded by Solar System
observations ($\dot{G}/G$ from LLR at $347\times$) through a result that is
not model-dependent at the DFD coupling level. Three resolution mechanisms
(environmental screening, modified evolution, observational model dependence)
have been assessed and found inadequate.

An additional internal inconsistency (T2) has been identified: the psi-field
equation is sourced only by EM stress-energy, yet the Mach integral that
generates $\psicosmo$ uses the identification $\psi = GM/c^2 r$ (the Newtonian
gravitational potential). These are inconsistent unless a matter source term
is added to the psi-field equation---a fundamental theoretical revision.

The P14+P5 incompatibility ($\psicosmo$ needed for MOND vs.\ $\psicosmo$
allowed by $\dot{G}/G$) is a derived consequence of T1 with a $60\times$
primary tension and an underlying $3.7\times$ tension between the Mach integral
and the MOND formula even without the observational bound.

The shadow $+4.6\%$ prediction emerges as the cleanest surviving P14 claim,
consistent with both M87* and Sgr A* EHT measurements at $< 1\sigma$, and
not dependent on the cosmological sector.

\textbf{T1 is a genuine partial falsification of DFD}: it falsifies the minimal
cosmological coupling sector, while leaving the field-theoretic core (mode spectrum,
dispersion relations, shadow mechanism, EM coupling) intact.

%=============================================================================
\end{document}
